\section{Background and Related
Work}\label{background-and-related-work}

\subsection{Data-Driven Business Process
Simulation}\label{data-driven-business-process-simulation}

Business process simulation estimates how a process may behave under
given assumptions about arrivals, routing, resources, calendars, service
times, and delays. In data-driven BPS, these assumptions are learned
from event logs rather than specified only by domain experts. Event logs
make it possible to estimate trace variant frequencies, activity
durations, case inter-arrival times, resource assignments, and workload
distributions. The simulator can then generate new event logs that are
compared with observed behavior or used for what-if analysis.

Early and mainstream data-driven BPS work focuses on discovering
executable simulation models from event logs. Camargo, Dumas, and
Gonzalez-Rojas~\cite{camargo2020automated} frame this as an
accuracy-optimized model discovery problem: instead of manually tuning simulation parameters, the
discovery procedure searches for configurations that maximize similarity
between the event log generated by a simulation model and the observed
log. This line of work is foundational for the present study because it
establishes the event log as both calibration data and evaluation
target. It also clarifies that simulation quality depends on multiple
learned components, including control-flow, arrival behavior, service
times, calendars, and resource assignments.

The difficulty is that high-quality BPS requires more than replaying
common traces. A simulator must reproduce multiple behavioral dimensions
simultaneously. A model that matches control-flow variants may still
fail on cycle time, resource workload, or circadian timing. Similarly, a
model that matches average cycle time may hide unrealistic resource
concentration or handover patterns. Chapela-Campa et al.~\cite{chapela2025quality} address
this problem by proposing a framework for measuring BPS model quality
across several log-distance dimensions. Their work is important for this
study because it provides an evaluation anchor that is stricter than
visual inspection or one-dimensional accuracy.

This evaluation perspective also motivates the conservative
interpretation of LLM-agent simulation. A generated log is not useful
merely because the simulated agents produce plausible explanations. It
must be compared with held-out behavior across control-flow, temporal,
and resource dimensions. The present study therefore uses Chapela-Campa
et al.'s framework not as an auxiliary check, but as the primary
safeguard against overclaiming the value of LLM-style agents.

\subsection{Resource Heterogeneity in
BPS}\label{resource-heterogeneity-in-bps}

Traditional BPS often attaches resources to a process model as
capacities, calendars, or pools. This representation is useful but can
flatten the heterogeneity of organizational work. In real processes,
resources differ in their skills, performance, availability,
interruption patterns, workload, and multitasking behavior. Several
recent BPS studies have therefore shifted from undifferentiated resource
pools toward individually parameterized resources.

Resource calendars are one example. Martin et al.~\cite{martin2020calendars} propose methods
for retrieving resource availability calendars from event logs, showing
how timestamps can be used to infer when resources are likely available.
Estrada-Torres et al.~\cite{estrada2020multitasking} extend the resource perspective by
addressing multitasking, a common source of mismatch between simulated
and real execution times. Their work shows that assuming monotasking
resources can overestimate execution times when real resources handle
overlapping work. Chapela-Campa and Dumas~\cite{chapela2024extraneous} further show that not
all waiting time can be explained by resource contention or
unavailability; some delays are extraneous to the modeled queue/resource
mechanism and need to be captured explicitly.

The strongest resource-centric bridge for this study is the work on
differentiated resources. Lopez-Pintado, Dumas, and Berx~\cite{lopezpintado2024differentiated} show
that treating each resource as an individual entity with its own
performance and availability model can improve the reproduction of
cycle-time distributions and work rhythm compared with undifferentiated
pools. Lopez-Pintado and Dumas~\cite{lopezpintado2025probabilistic} further model intermittent
availability and multitasking probabilistically, estimating availability
and multitasking probabilities from event logs. Together, these studies
show that many agent-like characteristics can already be derived
statistically from logs: capabilities, processing-time distributions,
calendars, intermittent availability, multitasking capacity, and delay
distributions.

This literature is directly relevant to the proposed artifact. It
suggests that an LLM-agent simulator should not start from open-ended
prompts. It should first inherit the strongest available data-driven
resource models. In the present implementation, this insight is
operationalized through log-derived capabilities, service-time samples,
waiting-time samples, case-arrival samples, and handover priors. The
LLM-agent proxy is therefore added on top of a resource-centric
statistical foundation rather than replacing it.

\subsection{Agent System Mining and Agent-Based Process
Simulation}\label{agent-system-mining-and-agent-based-process-simulation}

Resource-centric simulation reverses part of the process-centric
perspective. It asks what behavior can be learned about the actors
themselves: which activities they perform, how quickly they perform
them, how they interact, and how their local choices shape the global
process. Agent-based process simulation develops this idea further by
treating resources as autonomous agents rather than passive execution
capacity.

Agent System Mining (ASM) provides the conceptual link between process
mining and multi-agent simulation. Tour, Polyvyanyy, and Kalenkova
\cite{tour2021asm} argue that organizations can be viewed as multi-agent systems
whose behavior is recorded in event data. Instead of discovering only a
control-flow process model, ASM seeks to discover agents and
interactions among agents. Tour et al.~\cite{tour2023agentminer} instantiate this idea in
Agent Miner, an algorithm for discovering agent models and interaction
patterns from event data.

AgentSimulator is the most direct simulation starting point for this
project. It proposes a resource-first approach to data-driven BPS in
which agents correspond primarily to individual resources with
discovered behavior models \cite{kirchdorfer2024agentsimulator}. AgentSimulator is
important because it demonstrates that agent-based data-driven BPS can
be constructed from event logs and evaluated against conventional
process simulation approaches. The present study builds on the same
broad motivation but focuses on a narrower question: how an LLM-style
decision module can be inserted into a resource-centric simulator while
remaining constrained by log-derived behavior.

The relationship to ASM and AgentSimulator also clarifies the novelty
boundary. This study does not claim to introduce agent-based BPS from
scratch. Rather, it asks how the agent decision layer can be augmented
with an LLM-style local behavior module while retaining event-log
compatibility, action masks, statistical priors, and measurable BPS
quality. In this sense, the proposed artifact is positioned between
resource-centric BPS and LLM-based multi-agent simulation.

\subsection{LLM-Empowered Agent-Based Modeling and Agentic
BPM}\label{llm-empowered-agent-based-modeling-and-agentic-bpm}

LLM-based agents have become a prominent design pattern in multi-agent
systems. Compared with simple rule-based agents, LLM agents can produce
natural-language rationales, condition decisions on richer context, and
simulate heterogeneous behavioral styles. Gao et al.~\cite{gao2024llmabm} summarize
this opportunity for agent-based modeling and simulation, arguing that
LLMs can enrich agent cognition and interaction.

For business process simulation, these strengths are attractive but
insufficient by themselves. A process simulator must generate valid,
measurable process events. An LLM that produces plausible text but
invalid event-log behavior would weaken rather than strengthen BPS.
Therefore, the relevant research question is not whether an LLM can
describe a process, but whether LLM-like decision-making can be grounded
in event-log-derived constraints. This study uses the term LLM-agent
proxy for a local policy that occupies the same interface as a future
LLM call: it observes state, memory, feasible actions, resource priors,
and handover context, then selects an action and emits a reason. The
proxy allows the design to be tested without introducing API
nondeterminism, latency, cost, or prompt instability into the first
experiment.

The emerging literature on Agentic Business Process Management gives
this question a broader BPM context. Dumas, Milani, and Chapela-Campa
describe Agentic BPMS as a new class of platforms that combine
autonomy, reasoning, learning, and process mining. In such systems,
agents may sense process states, reason about improvement opportunities,
and act to maintain or optimize performance. This vision makes LLM-agent
process simulation relevant beyond a single simulator: simulation can
become a controlled environment for testing agent policies before they
are deployed in operational process management \cite{dumas2026agenticbpms}.

At the same time, Agentic BPM strengthens the need for guardrails. If
future process-aware agents are allowed to act autonomously, their
behavior must remain auditable and grounded in process evidence. The
proposed simulator contributes to this agenda by separating three
layers: a log-derived process environment, a constrained local decision
interface, and validation against held-out event logs. This separation
makes it possible to study LLM-style decision behavior without allowing
language generation to overwrite the empirical process structure.

\subsection{Simulation Verification and
Validation}\label{simulation-verification-and-validation}

Simulation research has long emphasized verification and validation.
Sargent~\cite{sargent2013vv} distinguishes whether a simulation is implemented
correctly from whether it is an adequate representation of the target
system. This distinction is essential for LLM-agent BPS. A system may
execute without errors and still be invalid as a process simulation.
Conversely, a failed metric can be diagnostically useful if it reveals
which behavioral assumption is missing. In this project, early temporal
failures led to the addition of empirical inter-event waiting times and
case inter-arrival sampling. The evaluation is therefore not only a
final scorecard; it is part of the artifact design cycle.

Taken together, the literature suggests a specific research gap.
Data-driven BPS provides log-derived control-flow, timing, and resource
parameters. Resource-centric BPS shows that individual resources can and
should be modeled heterogeneously. ASM and AgentSimulator show how event
logs can be interpreted as traces of interacting agents. LLM-empowered
ABM and Agentic BPM suggest richer reasoning and communication among
agents. What remains underdeveloped is an evaluable bridge between these
streams: a process-mining-compatible simulator in which an LLM-style
agent policy reasons locally, is constrained by log-derived feasible
actions, emits interpretable reasoning and handover traces, and is
evaluated with standard BPS log-distance measures.

